\IfFileExists{stacks-project.cls}{%
\documentclass{stacks-project}
}{%
\documentclass{amsart}
}
\usepackage{amssymb}

% For dealing with references we use the comment environment
\usepackage{verbatim}
\usepackage{fancyvrb}
\usepackage{pmboxdraw}
\newenvironment{reference}{\comment}{\endcomment}
%\newenvironment{reference}{}{}
\newenvironment{slogan}{\comment}{\endcomment}
\newenvironment{history}{\comment}{\endcomment}

% For source figures retained as assets
\usepackage{graphicx}

% For commutative diagrams we use Xy-pic
\usepackage[all]{xy}

% We use 2cell for 2-commutative diagrams.
\xyoption{2cell}
\UseAllTwocells

% We use multicol for the list of chapters between chapters
\usepackage{multicol}

% This is generally recommended for better output
\usepackage{lmodern}
\usepackage[T1]{fontenc}
\usepackage{textalpha}

% For cross-file-references
% We prefer the snapshot tmp/xref/#2.aux,
% which latexmkrc replaces atomically, so
% that documents compiled simultaneously
% never read a half-written .aux file.
% The last argument keeps links to #2.pdf.
\usepackage{xr-hyper}
\newcommand{\maybeexternaldocument}[2]{%
  \IfFileExists{#2.aux}{%
    \externaldocument[#1]{#2}%
  }{%
    \PackageWarning{stack}{Missing #2.aux; external references with prefix #1 unavailable}%
  }%
}
\let\externaldocumentorig\externaldocument
\renewcommand{\externaldocument}[2][]{%
  \IfFileExists{tmp/xref/#2.aux}{%
    \externaldocumentorig[#1]{tmp/xref/#2}[#2.pdf]%
  }{%
  \IfFileExists{#2.aux}{%
    \externaldocumentorig[#1]{#2}%
  }{%
    \PackageWarning{stack}{Missing #2.aux; external references with prefix #1 unavailable}%
  }%
  }%
}

% Package for hypertext links:
\usepackage{hyperref}

% For any local file, say "hello.tex" you want to link to please
% use \externaldocument[hello-]{hello}
\externaldocument[introduction-]{introduction}
\externaldocument[conventions-]{conventions}
\externaldocument[sets-]{sets}
\externaldocument[categories-]{categories}
\externaldocument[topology-]{topology}
\externaldocument[sheaves-]{sheaves}
\externaldocument[sites-]{sites}
\externaldocument[stacks-]{stacks}
\externaldocument[fields-]{fields}
\externaldocument[algebra-]{algebra}
\externaldocument[brauer-]{brauer}
\externaldocument[homology-]{homology}
\externaldocument[derived-]{derived}
\externaldocument[simplicial-]{simplicial}
\externaldocument[more-algebra-]{more-algebra}
\externaldocument[smoothing-]{smoothing}
\externaldocument[modules-]{modules}
\externaldocument[sites-modules-]{sites-modules}
\externaldocument[injectives-]{injectives}
\externaldocument[cohomology-]{cohomology}
\externaldocument[sites-cohomology-]{sites-cohomology}
\externaldocument[dga-]{dga}
\externaldocument[dpa-]{dpa}
\externaldocument[sdga-]{sdga}
\externaldocument[hypercovering-]{hypercovering}
\externaldocument[schemes-]{schemes}
\externaldocument[constructions-]{constructions}
\externaldocument[properties-]{properties}
\externaldocument[morphisms-]{morphisms}
\externaldocument[coherent-]{coherent}
\externaldocument[divisors-]{divisors}
\externaldocument[limits-]{limits}
\externaldocument[varieties-]{varieties}
\externaldocument[topologies-]{topologies}
\externaldocument[descent-]{descent}
\externaldocument[perfect-]{perfect}
\externaldocument[more-morphisms-]{more-morphisms}
\externaldocument[flat-]{flat}
\externaldocument[groupoids-]{groupoids}
\externaldocument[more-groupoids-]{more-groupoids}
\externaldocument[etale-]{etale}
\externaldocument[chow-]{chow}
\externaldocument[intersection-]{intersection}
\externaldocument[pic-]{pic}
\externaldocument[weil-]{weil}
\externaldocument[adequate-]{adequate}
\externaldocument[dualizing-]{dualizing}
\externaldocument[duality-]{duality}
\externaldocument[discriminant-]{discriminant}
\externaldocument[derham-]{derham}
\externaldocument[local-cohomology-]{local-cohomology}
\externaldocument[algebraization-]{algebraization}
\externaldocument[curves-]{curves}
\externaldocument[resolve-]{resolve}
\externaldocument[models-]{models}
\externaldocument[functors-]{functors}
\externaldocument[equiv-]{equiv}
\externaldocument[pione-]{pione}
\externaldocument[etale-cohomology-]{etale-cohomology}
\externaldocument[proetale-]{proetale}
\externaldocument[relative-cycles-]{relative-cycles}
\externaldocument[more-etale-]{more-etale}
\externaldocument[trace-]{trace}
\externaldocument[crystalline-]{crystalline}
\externaldocument[spaces-]{spaces}
\externaldocument[spaces-properties-]{spaces-properties}
\externaldocument[spaces-morphisms-]{spaces-morphisms}
\externaldocument[decent-spaces-]{decent-spaces}
\externaldocument[spaces-cohomology-]{spaces-cohomology}
\externaldocument[spaces-limits-]{spaces-limits}
\externaldocument[spaces-divisors-]{spaces-divisors}
\externaldocument[spaces-over-fields-]{spaces-over-fields}
\externaldocument[spaces-topologies-]{spaces-topologies}
\externaldocument[spaces-descent-]{spaces-descent}
\externaldocument[spaces-perfect-]{spaces-perfect}
\externaldocument[spaces-more-morphisms-]{spaces-more-morphisms}
\externaldocument[spaces-flat-]{spaces-flat}
\externaldocument[spaces-groupoids-]{spaces-groupoids}
\externaldocument[spaces-more-groupoids-]{spaces-more-groupoids}
\externaldocument[bootstrap-]{bootstrap}
\externaldocument[spaces-pushouts-]{spaces-pushouts}
\externaldocument[spaces-chow-]{spaces-chow}
\externaldocument[groupoids-quotients-]{groupoids-quotients}
\externaldocument[spaces-more-cohomology-]{spaces-more-cohomology}
\externaldocument[spaces-simplicial-]{spaces-simplicial}
\externaldocument[spaces-duality-]{spaces-duality}
\externaldocument[formal-spaces-]{formal-spaces}
\externaldocument[restricted-]{restricted}
\externaldocument[spaces-resolve-]{spaces-resolve}
\externaldocument[formal-defos-]{formal-defos}
\externaldocument[defos-]{defos}
\externaldocument[cotangent-]{cotangent}
\externaldocument[examples-defos-]{examples-defos}
\externaldocument[algebraic-]{algebraic}
\externaldocument[examples-stacks-]{examples-stacks}
\externaldocument[stacks-sheaves-]{stacks-sheaves}
\externaldocument[criteria-]{criteria}
\externaldocument[artin-]{artin}
\externaldocument[quot-]{quot}
\externaldocument[stacks-properties-]{stacks-properties}
\externaldocument[stacks-morphisms-]{stacks-morphisms}
\externaldocument[stacks-limits-]{stacks-limits}
\externaldocument[stacks-cohomology-]{stacks-cohomology}
\externaldocument[stacks-perfect-]{stacks-perfect}
\externaldocument[stacks-introduction-]{stacks-introduction}
\externaldocument[stacks-more-morphisms-]{stacks-more-morphisms}
\externaldocument[stacks-geometry-]{stacks-geometry}
\externaldocument[moduli-]{moduli}
\externaldocument[moduli-curves-]{moduli-curves}
\externaldocument[examples-]{examples}
\externaldocument[exercises-]{exercises}
\externaldocument[guide-]{guide}
\externaldocument[desirables-]{desirables}
\externaldocument[coding-]{coding}
\externaldocument[obsolete-]{obsolete}
\externaldocument[fdl-]{fdl}
\externaldocument[index-]{index}

% Heap Project documents
\externaldocument[k3-]{k3}
\externaldocument[torelli-]{torelli}
\externaldocument[linear-algebra-]{linear-algebra}
\externaldocument[groups-]{groups}
\externaldocument[complex-analysis-]{complex-analysis}
\externaldocument[functional-analysis-]{functional-analysis}
\externaldocument[categories-]{categories}
\externaldocument[commutative-algebra-]{commutative-algebra}
\externaldocument[ringed-spaces-]{ringed-spaces}
\externaldocument[homological-algebra-]{homological-algebra}
\externaldocument[lie-algebras-]{lie-algebras}
\externaldocument[representation-theory-]{representation-theory}
\externaldocument[smooth-manifolds-]{smooth-manifolds}
\externaldocument[differential-topology-]{differential-topology}
\externaldocument[algebraic-topology-]{algebraic-topology}
\externaldocument[differential-geometry-]{differential-geometry}
\externaldocument[algebraic-geometry-]{algebraic-geometry}
\externaldocument[symplectic-geometry-]{symplectic-geometry}
\externaldocument[complex-geometry-]{complex-geometry}
\externaldocument[hodge-theory-]{hodge-theory}
\externaldocument[%
variations-of-hodge-structures-]{%
variations-of-hodge-structures}
\externaldocument[%
homological-mirror-symmetry-]{%
homological-mirror-symmetry}
\maybeexternaldocument{deformations-}{deformations}
\externaldocument[vertex-algebras-]{vertex-algebras}
\externaldocument[springer-]{springer}
\externaldocument[infty-categories-]{infty-categories}
\externaldocument[seminars-differential-geometry-]{%
seminars-differential-geometry}
\externaldocument[seminars-complex-geometry-]{%
seminars-complex-geometry}
\externaldocument[seminars-algebraic-geometry-]{%
seminars-algebraic-geometry}
\externaldocument[seminars-analysis-pde-]{%
seminars-analysis-pde}
\externaldocument[surfaces-]{surfaces}
\externaldocument[birational-maps-commutative-algebra-]{birational-maps-commutative-algebra}
\externaldocument[cremona-transformations-]{cremona-transformations}
\externaldocument[derived-categories-of-sheaves-]{derived-categories-of-sheaves}
\externaldocument[geometric-invariant-theory-]{geometric-invariant-theory}
\externaldocument[geometric-stability-conditions-and-group-actions-]{geometric-stability-conditions-and-group-actions}
\externaldocument[moduli-spaces-of-sheaves-]{moduli-spaces-of-sheaves}
\externaldocument[bridgeland-stability-general-theory-]{bridgeland-stability-general-theory}
\externaldocument[hilbert-schemes-]{hilbert-schemes}
\externaldocument[noncommutative-abelian-surfaces-and-kummer-type-hyperkahler-varieties-]{noncommutative-abelian-surfaces-and-kummer-type-hyperkahler-varieties}
\externaldocument[mumford-tate-groups-in-hodge-theory-]{mumford-tate-groups-in-hodge-theory}
\externaldocument[hodge-birational-atoms-2026-]{hodge-birational-atoms-2026}
\externaldocument[xxii-egd-2026-]{%
xxii-egd-2026}
\externaldocument[pde-]{pde}
\externaldocument[probability-]{probability}

\externaldocument[linguistics-]{linguistics}

\externaldocument[%
16th-alga-meeting-2026-]{%
16th-alga-meeting-2026}

\externaldocument[%
iv-jornadas-lefschetz-2026-]{%
iv-jornadas-lefschetz-2026}

\externaldocument[reading-the-masters-]{%
reading-the-masters}

\externaldocument[real-analysis-]{%
real-analysis}

\externaldocument[optimal-transport-]{%
optimal-transport}

\externaldocument[history-]{history}

\externaldocument[seiberg-witten-]{seiberg-witten}

\externaldocument[history-philosophy-biology-]{history-philosophy-biology}

\externaldocument[classical-mechanics-]{classical-mechanics}
\externaldocument[electromagnetism-]{electromagnetism}
\externaldocument[waves-and-optics-]{waves-and-optics}
\externaldocument[wif-]{waves-in-fluids-iv}
\externaldocument[thermal-physics-]{thermal-physics}
\externaldocument[quantum-mechanics-]{quantum-mechanics}
\externaldocument[special-relativity-]{special-relativity}
\externaldocument[physics-]{physics}

\externaldocument[chemistry-]{chemistry}

\externaldocument[molecular-cell-biology-]{molecular-cell-biology}

\externaldocument[microbiology-]{microbiology}

\externaldocument[protistology-]{protistology}

\externaldocument[mycology-]{mycology}

\externaldocument[botany-]{botany}

\externaldocument[zoology-]{zoology}

\externaldocument[animal-histology-]{animal-histology}

\externaldocument[animal-physiology-]{animal-physiology}

\externaldocument[systematics-biogeography-]{systematics-biogeography}

\externaldocument[ecology-natural-resources-]{ecology-natural-resources}

\externaldocument[earth-sciences-]{earth-sciences}

\externaldocument[biostatistics-]{biostatistics}

\externaldocument[mathematical-systems-biology-]{mathematical-systems-biology}

\externaldocument[cimpa-el-salvador-2024-]{cimpa-el-salvador-2024}

\externaldocument[anthropology-]{anthropology}

\externaldocument[yang-mills-theory-]{yang-mills-theory}

\externaldocument[seminars-algebra-]{seminars-algebra}

\externaldocument[seminars-physics-]{seminars-physics}

\externaldocument[seminars-computing-]{seminars-computing}

\externaldocument[oktobermat-2026-]{oktobermat-2026}

\externaldocument[saet-2026-]{saet-2026}

\externaldocument[seminars-symplectic-geometry-]{seminars-symplectic-geometry}

\externaldocument[seminars-dynamics-]{seminars-dynamics}

\externaldocument[oktobermat-2024-]{oktobermat-2024}

\externaldocument[quantum-groups-]{quantum-groups}

\externaldocument[legal-2027-]{legal-2027}

\externaldocument[derived-categories-]{derived-categories}

% Theorem environments.
%
\theoremstyle{plain}
\newtheorem{theorem}[subsection]{Theorem}
\newtheorem{proposition}[subsection]{Proposition}
\newtheorem{lemma}[subsection]{Lemma}
\newtheorem{corollary}[subsection]{Corollary}

\theoremstyle{definition}
\newtheorem{definition}[subsection]{Definition}
\newtheorem{example}[subsection]{Example}
\newtheorem{exercise}[subsection]{Exercise}
\newtheorem{situation}[subsection]{Situation}

\theoremstyle{remark}
\newtheorem{remark}[subsection]{Remark}
\newtheorem{remarks}[subsection]{Remarks}

\numberwithin{equation}{subsection}

% Macros
%
\def\lim{\mathop{\mathrm{lim}}\nolimits}
\def\colim{\mathop{\mathrm{colim}}\nolimits}
\def\Spec{\mathop{\mathrm{Spec}}}
\def\Hom{\mathop{\mathrm{Hom}}\nolimits}
\def\Ext{\mathop{\mathrm{Ext}}\nolimits}
\def\SheafHom{\mathop{\mathcal{H}\!\mathit{om}}\nolimits}
\def\SheafExt{\mathop{\mathcal{E}\!\mathit{xt}}\nolimits}
\def\Sch{\mathit{Sch}}
\def\Mor{\mathop{\mathrm{Mor}}\nolimits}
\def\Ob{\mathop{\mathrm{Ob}}\nolimits}
\def\Sh{\mathop{\mathit{Sh}}\nolimits}
\def\NL{\mathop{N\!L}\nolimits}
\def\CH{\mathop{\mathrm{CH}}\nolimits}
\def\proetale{{pro\text{-}\acute{e}tale}}
\def\etale{{\acute{e}tale}}
\def\QCoh{\mathit{QCoh}}
\def\Ker{\mathop{\mathrm{Ker}}}
\def\Im{\mathop{\mathrm{Im}}}
\def\Coker{\mathop{\mathrm{Coker}}}
\def\Coim{\mathop{\mathrm{Coim}}}

% Boxtimes
%
\DeclareMathSymbol{\boxtimes}{\mathbin}{AMSa}{"02}

%
% Macros for moduli stacks/spaces
%
\def\QCohstack{\mathcal{QC}\!\mathit{oh}}
\def\Cohstack{\mathcal{C}\!\mathit{oh}}
\def\Spacesstack{\mathcal{S}\!\mathit{paces}}
\def\Quotfunctor{\mathrm{Quot}}
\def\Hilbfunctor{\mathrm{Hilb}}
\def\Curvesstack{\mathcal{C}\!\mathit{urves}}
\def\Polarizedstack{\mathcal{P}\!\mathit{olarized}}
\def\Complexesstack{\mathcal{C}\!\mathit{omplexes}}
% \Pic is the operator that assigns to X its picard group, usage \Pic(X)
% \Picardstack_{X/B} denotes the Picard stack of X over B
% \Picardfunctor_{X/B} denotes the Picard functor of X over B
\def\Pic{\mathop{\mathrm{Pic}}\nolimits}
\def\Picardstack{\mathcal{P}\!\mathit{ic}}
\def\Picardfunctor{\mathrm{Pic}}
\def\Deformationcategory{\mathcal{D}\!\mathit{ef}}


\begin{document}

\title{Yang-Mills theory}
\maketitle

\phantomsection
\label{section-phantom}

\tableofcontents

\section{About these notes}

These are reading notes on the Yang--Mills
existence and mass gap problem, followed by
notes from a seminar. They are not a survey.
The sources are Jaffe and Witten
\cite{yangmills}, Douglas
\cite{reportdouglas} and a video lecture by
Gross \cite{gross-video}. I have not checked
every statement against these sources, and I
mark what I have not understood.

\section{Maxwell's theory as a gauge theory}

Electromagnetism is the basic example of a
gauge theory. Its gauge group is
$\mathrm{U}(1)$. The dynamical field is a
connection $A$ on a $\mathrm{U}(1)$-bundle
over Minkowski space $\mathbb{R}^{3,1}$,
which is locally a 1-form. The observable is
the field strength, that is, the curvature
$$
F=dA,
$$
whose components are the familiar electric
and magnetic fields. In vacuum Maxwell's
equations read
$$
dF=0,\qquad d{*}F=0,
$$
where $*$ is the Hodge star. They are second
order PDEs for $A$, and they are the
Euler--Lagrange equations of the action
$$
S=\frac{1}{2e^2}\int F\wedge {*F}
$$
(up to sign and normalization conventions).
The constant $e$ is the electric charge. The
action is invariant under the Lorentz group
$\mathrm{SO}(3,1)$. Light is a solution: a
ripple of the field strength travelling at
the speed of light.

In physics one chooses units and keeps track
of universal constants, such as the speed of
light $c$ and Planck's constant $\hbar$. If
we set $\hbar=c=1$, every quantity becomes a
power of a length, and the charge $e$ becomes
a dimensionless number. With a suitable
normalization, $e^2/\hbar c$ is the
fine-structure constant, which is about
$1/137$.

Electrodynamics in vacuum is a free theory,
and it is easy to quantize: one can construct
the Hilbert space explicitly. Nobody has
proved the existence of quantum
electrodynamics as a theory in the continuum
(Maxwell's field coupled to electrons), as I
understand from Gross's lecture below.

\section{Yang--Mills theory}

Yang--Mills theory, also called non-abelian
gauge theory, replaces $\mathrm{U}(1)$ by a
compact gauge group $G$. The curvature of a
connection $A$ becomes
$$
F=dA+A\wedge A,
$$
and Maxwell's equations are replaced by the
Yang--Mills equations
$$
d_AF=0,\qquad d_A{*}F=0,
$$
where $d_A$ is the gauge-covariant extension
of the exterior derivative. I have not yet
worked out how the Hodge star acts here.

The theory comes in a classical version,
which is a system of PDEs, and a quantum
version, which has no satisfactory
mathematical definition in dimension four.
Jaffe and Witten \cite{yangmills} recall two
historical obstacles to applying the theory
to the nuclear forces. Classical Yang--Mills
waves are massless. For the weak force this
was overcome in the 1960s and 1970s by
introducing an additional Higgs field. For
the strong force the solution is of a
different nature: asymptotic freedom, the
fact that at short distances the quantum
field behaves much like the classical one,
while at long distances the classical theory
is no longer a good guide.

\section{The Millennium problem}

The problem, as stated by Jaffe and Witten
\cite{yangmills}, is to prove that for any
compact simple gauge group $G$, a non-trivial
quantum Yang--Mills theory exists on
$\mathbb{R}^4$ and has a mass gap $\Delta>0$.
Here existence includes establishing
axiomatic properties for the theory, which
are listed in their paper. The mass gap says
that the quantum particles have positive
masses, even though the classical waves
travel at the speed of light.

Douglas \cite{reportdouglas} explains the
problem through a lattice. Take a lattice
$\Gamma$ such as $\mathbb{Z}^4\subset
\mathbb{R}^4$ and give it a graph structure.
A $G$-connection is a function from the edges
to $G$. One defines its curvature, and the
Yang--Mills action is the sum of the squares
of the curvatures. To quantize, one defines
on a finite subgraph $\gamma\subset\Gamma$ a
functional integral $Z$, called the
\emph{partition function}. It is of physical
interest, as are the expectation values of
quantum observables. The objective is to
define a functional integral on all of
$\Gamma$ by taking the limit of $Z$ over
larger and larger $\gamma$. Existence and the
mass gap are properties that this
construction should have. There are other,
very different ways to think about the
theory.

Douglas also mentions two families of quantum
field theories believed to be similar to
four-dimensional Yang--Mills theory: the
two-dimensional nonlinear sigma models with
target a group manifold, and supersymmetric
Yang--Mills theories, where one also has
fermions. There is also the 1994 work
of Seiberg and Witten
(see the chapter on
Seiberg--Witten theory), and the
integrability of $N=4$ supersymmetric
Yang--Mills theory, which is related to the
AdS/CFT correspondence. I have not checked
these two items.

\section{Gross's lecture: why this problem?}

In his video lecture \cite{gross-video},
Gross gives two reasons for the prize. First,
Yang--Mills theories are the standard theory
of elementary particles, the most complete
and ambitious physical theory we have.
Second, physicists have ways to discover
truth that are not guided solely by
mathematics. It is hard to say how a proof
could even fail, since we do not know what it
would mean to prove the non-existence of a
quantum field theory. One could define a
Yang--Mills theory without a positive mass
gap, but that would shock everybody, because
nature, experiments and analogies tell us the
theory is correct. It is like walking on a
rope over a precipice: nature guides us.
Gross takes strong bets against the
non-existence of the mass gap. There is no
mathematics yet that explains this.

Quantum field theory assumes that fields are
functions on spacetime, that is, that space
is a continuum which can be probed at
arbitrarily short distances. If one puts a
cutoff at some very short distance, the
theory of quantum electrodynamics is well
defined, and asymptotic expansions give good
approximations. But the theory probably does
not exist down to arbitrarily short
distances. For the same reasons, Gross says,
non-abelian gauge theories, even in vacuum,
do exist at arbitrarily short distances: they
are complete theories.

\section{What is a quantum field theory?}

Gross's question is: what does it mean to
define a quantum field theory, or to prove
that it exists? We need the following.

\begin{enumerate}
\item A Hilbert space $\mathcal{H}$ of states
$|\psi\rangle$. The norm of a state gives a
probability, so it must be positive.
\item Relativistic invariance: Poincar\'e
symmetry, consisting of translations
$x\mapsto x+a$ in space and time and Lorentz
transformations $\mathrm{SO}(3,1)$. In
quantum mechanics these symmetries are
represented by unitary operators, under which
the states and the fields transform
covariantly.
\item A vacuum state $|\Omega\rangle$, the
ground state of the system, invariant under
translations and Lorentz transformations.
\item Spectral properties for the generators
of the translations, the momenta $P^\mu$. In
particular the energy $P^0$ has spectrum
bounded below, $P^0\geq 0$. The mass is
defined by $m^2=E^2-|P|^2$.
\end{enumerate}

In electromagnetism the spectrum of the
energy is continuous. In a theory with a mass
gap, the spectrum has the vacuum
$|\Omega\rangle$, then a discrete state of
mass $m_G$, where $G$ stands for gap, and
then a continuum of multiparticle states,
starting at $2m_G$ plus some kinetic energy,
perhaps with some isolated states in between.
So there are no massless particles like light
rays: there is a gap between the vacuum and
the first state, whose energy is the mass of
the lightest particle. The interesting
dynamical issue is how to produce, from a
theory with the same structure as classical
electromagnetism, a mass gap. Quantum effects
produce a scale, the mass gap, while the
classical theory has none.

\section{Seminar talk on quantum fields}

These are notes from a seminar talk, whose
speaker may share better notes. The plan was:
quantum field theory, axiomatic quantum field
theory, and quantum Yang--Mills theory, with
the mass gap, perturbative versus lattice
methods, the case of dimension $2$, and the
difficulties.

The Standard Model has two kinds of
particles. The fermions are quarks ($u$, $c$,
$t$, $d$, $s$, $b$) and leptons (the electron
$e^-$, $\mu^-$, $\tau^-$ and the neutrinos
$\nu_e$, $\nu_\mu$, $\nu_\tau$). The bosons
carry the forces: the photon $\gamma$, the
particles $Z$ and $W$ responsible for some
radioactive decays, the gluons $g$, and the
scalar Higgs boson $H$. Each force has a
gauge group: electromagnetism (and its
quantum version, quantum electrodynamics) has
$\mathrm{U}(1)$, with the photon; the
electroweak force has
$\mathrm{SU}(2)\times\mathrm{U}(1)$, and
involves the Higgs boson; quantum
chromodynamics has $\mathrm{SU}(3)$, with the
gluon. A proton is made of three quarks,
$uud$, and a neutron is $udd$, so they are
not fundamental. Almost every particle has an
antiparticle with the same mass and opposite
charge.

In the classical picture, the theory is
``deterministic'': there is a space of fields
$A$, an action functional
$S(A)=\int\mathcal{L}(A)$, and the equations
of motion are the Euler--Lagrange equations
$\delta S=0$. In the quantum picture, the
theory is ``probabilistic'': a particle
becomes a wave, as in the double-slit
experiment. One computes the Euclidean
partition function and expectation values of
observables $O$,
$$
Z=\int e^{-S(A)/\hbar}\,dA,\qquad
\langle O\rangle=
\frac{1}{Z}\int O(A)\,e^{-S(A)/\hbar}\,dA.
$$
The measure $dA$ is usually ill defined, but
there are ways around this in some cases. As
$\hbar\to 0$ the integral localizes at the
minima of $S$, which are solutions of the
Euler--Lagrange equations.

Here are examples of such theories.

\begin{enumerate}
\item On a $3$-manifold $M$ with a compact
gauge group $G$, the Chern--Simons theory has
action
$$
S(A)=\frac{k}{4\pi}\int_M
\mathrm{tr}\left(A\wedge dA+
\frac{2}{3}A\wedge A\wedge A\right),
$$
and partition function $Z=\int
e^{iS(A)/\hbar}\,dA$.
\item If the fields are the Riemannian
metrics on $M$, one gets quantum gravity.
\item Fix an auxiliary manifold $N$ and take
as fields the functions $M\to N$. This is a
sigma model. The sigma model in dimension one
is quantum mechanics.
\item Fix a principal $G$-bundle $P\to M$.
The fields are $G$-connections, the gauge
fields, and sections of associated vector
bundles, the matter fields. This is quantum
gauge theory.
\end{enumerate}

\section{Topological quantum field theory}

The same seminar series moved to topological
quantum field theory. In Atiyah's
formulation, a TQFT in dimension $n$ assigns
to every closed oriented $(n-1)$-manifold $Y$
a vector space $\mathcal{H}(Y)$, and to every
compact $n$-manifold $X$ with boundary
$\partial X=Y$ a vector
$Z(X)\in\mathcal{H}(Y)$. The axioms include
$$
\mathcal{H}(Y\sqcup Y')=
\mathcal{H}(Y)\otimes\mathcal{H}(Y'),
\qquad
\mathcal{H}(\overline{Y})=\mathcal{H}(Y)^*,
$$
together with a gluing axiom, which says that
gluing manifolds along a common boundary
component contracts the corresponding
vectors. In a $1$-extended theory, one also
assigns a category to every closed surface.

\bibliography{my}
\bibliographystyle{amsalpha}

\end{document}
